\documentclass[11pt]{article}
\usepackage[margin=2.2cm]{geometry}
\usepackage{booktabs}
\usepackage{amsmath,amssymb}
\usepackage[most]{tcolorbox}
\usepackage{xcolor}
\usepackage{tikz}
\usepackage{enumitem}

\newtcolorbox{concept}{colback=gray!5!white, colframe=gray!60!black,
  fonttitle=\bfseries, boxrule=0.5pt, arc=2pt}

\title{\textbf{How It Works --- Key Concepts}\\[4pt]
\large Neural Polar Decoding for Two-User MAC}
\author{}
\date{April 2026}

\begin{document}
\maketitle
\thispagestyle{empty}

%% ============================================================
\begin{concept}[title=1. MAC Polar Codes]
Two users $U$ and $V$ each independently encode their message bits using a standard polar encoder ($G_N = B_N F^{\otimes n}$: bit-reversal then XOR butterfly).
Both transmit simultaneously; the receiver sees $Z = f(X, Y) + \text{noise}$.
A \textbf{path} $\pi \in \{0,1\}^{2N}$ determines the decoding order: at each of $2N$ leaf positions, the decoder decides one bit of either $U$ or $V$.
Class B (symmetric rate) uses an interleaved path; Class C (corner rate) decodes all of $U$ first, then $V$.
\end{concept}

%% ============================================================
\begin{concept}[title=2. The Two Core Operations]
All polar decoding reduces to two tensor operations on $2 \times 2$ probability tables (one dimension per user):

\medskip
\textbf{Circular convolution (\texttt{circ\_conv}):} Marginalizes over an unknown bit linked by XOR encoding.
Given a tensor $T(a, b)$ representing $\Pr[\text{data} \mid u_{\text{left}}=a, v_{\text{right}}=b]$, convolution sums over the values consistent with $u_{\text{left}} \oplus u_{\text{right}} = c$:
\[
  (\texttt{circ\_conv}\; T)[c] = \sum_{a \oplus b = c} T(a, b)
\]
This is exactly circular convolution of each row/column. Used in \texttt{CalcLeft}.

\medskip
\textbf{Normalized product (\texttt{norm\_prod}):} Combines two independent pieces of evidence about the \emph{same} variable. If $T_1$ and $T_2$ are both likelihoods for the same $(u, v)$:
\[
  (\texttt{norm\_prod}(T_1, T_2))[a, b] \propto T_1(a, b) \cdot T_2(a, b)
\]
This is pointwise multiplication followed by normalization. Used in \texttt{CalcRight} and at leaves.
\end{concept}

%% ============================================================
\begin{concept}[title=3. CalcLeft / CalcRight / CalcParent]
The decoder operates on a binary tree with $N$ leaves. Each edge stores a $2 \times 2$ tensor.

\begin{itemize}[nosep]
  \item \textbf{CalcLeft} (top-down): Given the parent's top-down message, compute the left child's message using \texttt{circ\_conv}. This marginalizes over the right child's unknown contribution.
  \item \textbf{CalcRight} (top-down): Given the parent's top-down message and the left child's decision, compute the right child's message using \texttt{norm\_prod}. The left child's value is now known, so we condition on it.
  \item \textbf{CalcParent} (bottom-up): After deciding a leaf, propagate the decision upward. Converts a child's ``as-child'' tensor to the parent's ``as-parent'' tensor using \texttt{circ\_conv}.
\end{itemize}

These three operations are applied recursively, navigating the tree in bit-reversal order to decode all $2N$ bits.
\end{concept}

%% ============================================================
\begin{concept}[title=4. Chained Decoding (Two-Stage)]
For the MAC, decoding is \textbf{chained}:
\begin{enumerate}[nosep]
  \item \textbf{Stage 1:} Decode user $U$, treating $V$ as noise. The channel likelihoods marginalize over all possible $V$ values.
  \item \textbf{Stage 2:} Given the decoded $\hat{U}$, subtract its contribution from the received signal and decode $V$.
\end{enumerate}
For Class C (corner rate), the path naturally separates the two stages.
For Class B (interleaved path), $U$ and $V$ bits are decoded alternately --- each decision affects both users.
\end{concept}

%% ============================================================
\begin{concept}[title=5. The Neural Replacement (NPD / NCG)]
Replace exact $2 \times 2$ tensors with learned $d$-dimensional embeddings:
\begin{itemize}[nosep]
  \item Each edge stores a vector $\mathbf{e} \in \mathbb{R}^d$ instead of a $2 \times 2$ table.
  \item \texttt{CalcLeft}, \texttt{CalcRight}, \texttt{CalcParent} become \textbf{learned MLPs}: $f_\theta: \mathbb{R}^d \to \mathbb{R}^d$.
  \item At each leaf, a \textbf{decision head} maps the $d$-dim embedding to a bit probability.
  \item Channel observations $z$ are mapped to initial $d$-dim embeddings via a learned \textbf{z\_encoder}.
\end{itemize}

\textbf{NPD} (Neural Polar Decoder): replaces CalcLeft/Right/Parent with neural networks.\\
\textbf{NCG} (Neural Computational Graph): same idea applied to the full SC computational graph, using the interleaved tree structure from Ren et al.\ (2025).

Key advantage: the neural operations can implicitly learn channel structure (noise distribution, memory, interference) without explicit channel modeling.
\end{concept}

%% ============================================================
\begin{concept}[title=6. Training: \texttt{fast\_ce}]
Training is \textbf{teacher-forced and parallel}:
\begin{itemize}[nosep]
  \item Generate random codewords and channel outputs.
  \item Run the full tree computation with \emph{ground-truth} bits at each decision (teacher forcing). No sequential error propagation during training.
  \item Cross-entropy loss at every leaf: $\mathcal{L} = -\sum_{t=1}^{2N} \log p_\theta(\hat{b}_t = b_t)$.
  \item All $2N$ leaves computed in one forward pass; gradient depth is $O(\log N)$ (tree depth), not $O(N)$.
\end{itemize}
This enables training at $N = 1024$ in reasonable time. At \emph{test} time, decoding is sequential (each decision feeds into the next).
\end{concept}

%% ============================================================
\begin{concept}[title=7. Three-Phase Design]
\begin{enumerate}[nosep]
  \item \textbf{Rate-1 training:} Train the neural decoder at rate 1 (no frozen bits). The model learns the channel and tree operations without needing a good frozen-set design.
  \item \textbf{MI ranking:} Use the trained model's leaf cross-entropies as mutual information estimates. Rank bit positions by reliability. Select the $k_u + k_v$ most reliable positions as information bits.
  \item \textbf{Retrain at target rate:} Freeze the selected positions and retrain the decoder. The frozen-set and decoder are now \emph{co-adapted}.
\end{enumerate}
This avoids the chicken-and-egg problem: you need a good frozen set to train a decoder, but you need a decoder to evaluate frozen-set quality.
\end{concept}

\end{document}
